% !TeX root = semana03.tex
% !TeX program = pdflatex
% Material de la semana 03, ajustado al avance real del curso.
\documentclass[10pt,a5paper]{article}
\newif\ifmostrarSoluciones
\mostrarSolucionestrue
\usepackage{tikz}
\usepackage{estilo_apuntes}
\usetikzlibrary{arrows.meta,positioning,babel}
\encabezado{Semana 3}{Apunte semanal}

\begin{document}
\tituloclase{Límites: aproximación y cálculo}{Límites básicos, comparación y teorema del Sanduche}{Semana 3 · 5 horas}

\begin{objetivos}
\begin{itemize}
  \item Interpretar límites mediante tablas y gráficas.
  \item Distinguir el límite del valor de la función y reconocer su unicidad.
  \item Utilizar funciones localmente iguales y el álgebra de límites.
  \item Calcular límites polinómicos, racionales y radicales mediante sustitución, factorización y racionalización.
  \item Aplicar comparación, preservación del signo y el teorema del Sanduche a funciones acotadas y oscilatorias.
\end{itemize}
\end{objetivos}

\begin{resumen}
El límite describe el comportamiento de una función cerca de un punto, sin depender de su valor en ese punto. Primero exploramos ese comportamiento; después lo justificamos mediante reglas algebraicas y comparaciones. 
\end{resumen}

\section*{Aproximación a un punto}

\begin{definicion}[noción intuitiva de límite]
Escribimos
\[
\lim_{x\to a}f(x)=L
\]
cuando los valores $f(x)$ pueden hacerse tan próximos a $L$ como se desee al tomar $x$ suficientemente próximo a $a$, con $x\neq a$. El límite estudia valores en un entorno alrededor de $a$ del que se excluye el propio punto $a$.
\end{definicion}

\begin{nota}
La expresión $x\to a$ no significa que $x$ llegue a ser igual a $a$. Las afirmaciones
\[
  \lim_{x\to a}f(x)=L
  \qquad\text{y}\qquad
  f(a)=L
\]
son diferentes y ninguna implica por sí sola la otra.
\end{nota}

\begin{ejemplo}[exploración mediante una tabla]
Considere la función $f\colon\mathbb{R}\setminus\{1\}\to\mathbb{R}$ definida por
\[
f(x)=\frac{x^2-1}{x-1}.
\]
Complete o verifique la tabla y conjeture el límite cuando $x\to1$.
\[
\begin{array}{c|cccccc}
x&0.9&0.99&0.999&1.001&1.01&1.1\\ \hline
f(x)&1.9&1.99&1.999&2.001&2.01&2.1
\end{array}
\]
\end{ejemplo}
\ifmostrarSoluciones
\begin{solucion}
Los valores se aproximan a $2$ desde ambos lados. Como $x^2-1=(x-1)(x+1)$,
\[
f(x)=x+1\quad(x\neq1),
\]
y la evidencia sugiere
\[
  \lim_{x\to1}f(x)=2,
\]
aunque $f(1)$ no está definida.
\end{solucion}
\else\vspace{12mm}\fi

\begin{center}
\begin{tikzpicture}[x=1.05cm,y=0.8cm,every node/.style={font=\scriptsize}]
  \draw[Gris!70] (-0.4,0)--(3.1,0) node[right] {$x$};
  \draw[Gris!70] (0,-0.3)--(0,3.4) node[above] {$y$};
  \draw[Azul,very thick,domain=-0.1:2.8,samples=2] plot (\x,{\x+1});
  \draw[fill=white,draw=Rojo,thick] (1,2) circle (2.2pt);
  \draw[Gris,dashed] (1,0)--(1,2);
  \node[below] at (1,0) {$1$};
  \node[left] at (0,2) {$2$};
  \node[right] at (1,2) {valor al que se aproxima};
\end{tikzpicture}
\end{center}

\begin{definicion}[límites laterales]
El \textbf{límite por la izquierda}
\[
  \lim_{x\to a^-}f(x)
\]
utiliza valores $x<a$. El \textbf{límite por la derecha}
\[
  \lim_{x\to a^+}f(x)
\]
utiliza valores $x>a$.
\end{definicion}

\begin{proposicion}[límites laterales]
El límite existe y vale $L$ si y solo si existen los dos límites laterales y coinciden:
\[
\lim_{x\to a}f(x)=L
\iff
\lim_{x\to a^-}f(x)=L=\lim_{x\to a^+}f(x).
\]
Si los límites laterales son distintos, el límite no existe.
\end{proposicion}

\begin{ejemplo}[salto]
Considere la función $s:\mathbb{R}\to\mathbb{R}$ definida por
\[
s(x)=
\begin{cases}
  -1, & x<0,\\
  2, & x\geq0.
\end{cases}
\]
Determine los límites laterales de $s$ cuando $x$ se aproxima a $0$ y decida si existe
\[
  \lim_{x\to0}s(x).
\]
\end{ejemplo}
\ifmostrarSoluciones
\begin{solucion}
\[
\lim_{x\to0^-}s(x)=-1,\qquad
\lim_{x\to0^+}s(x)=2.
\]
Como son distintos,
\[
  \lim_{x\to0}s(x)
\]
no existe, aunque $s(0)=2$ sí está definido.
\end{solucion}
\else\vspace{13mm}\fi

\section*{Unicidad y valor de la función en el punto}

\begin{teorema}[unicidad del límite]
Si
\[
  \lim_{x\to a}f(x)
\]
existe, su valor es único. En consecuencia, una misma función no puede aproximarse simultáneamente a dos números distintos cuando $x\to a$.
\end{teorema}

\begin{ejercicios}[noción intuitiva de límite]
\begin{enumerate}
  \item Construya una tabla para estimar $\displaystyle\lim_{x\to2}(x^2+1)$.
  \item Estime $\displaystyle\lim_{x\to3}\dfrac{x^2-9}{x-3}$ con tres valores a cada lado de $3$.
  \item Dibuje una función con $f(2)=4$ y $\displaystyle\lim_{x\to2}f(x)=1$.
  \item Sea $f\colon\mathbb{R}\setminus\{0\}\to\mathbb{R}$, definida por $f(x)=|x|/x$. Explique si existe el límite de $f$ en $0$.
\end{enumerate}
\end{ejercicios}

\ifmostrarSoluciones
\begin{solucion}[de los ejercicios]
\begin{enumerate}
  \item Los valores se aproximan a $5$.
  \item Para $x\neq3$, el cociente es $x+3$; la tabla se aproxima a $6$.
  \item Por ejemplo, $f\colon\mathbb{R}\to\mathbb{R}$, definida por
  \[
    f(x)=
    \begin{cases}
      1, & x\neq2,\\
      4, & x=2.
    \end{cases}
  \]
  \item No existe: los límites laterales son $-1$ y $1$.
\end{enumerate}
\end{solucion}
\fi
\section*{Funciones localmente iguales}

\begin{proposicion}[igualdad en un entorno]
Suponga que existe un intervalo abierto $I$ que contiene a $a$, tal que $f$ y $g$ están definidas en $I\setminus\{a\}$ y satisfacen
\[
  f(x)=g(x)\qquad\text{para todo }x\in I\setminus\{a\}.
\]
Entonces
\[
  \lim_{x\to a}f(x)
  \quad\text{existe si y solo si}\quad
  \lim_{x\to a}g(x)
  \quad\text{existe}.
\]
En tal caso,
\[
  \lim_{x\to a}f(x)=\lim_{x\to a}g(x).
\]
Los valores $f(a)$ y $g(a)$ no intervienen: pueden ser distintos o incluso no estar definidos.
\end{proposicion}

\begin{ejemplo}[el valor aislado no determina el límite]
Sean $g,h\colon\mathbb{R}\to\mathbb{R}$, definidas por $g(x)=x+1$ y
\[
  h(x)=
  \begin{cases}
    x+1, & x\neq1,\\
    7, & x=1.
  \end{cases}
\]
Compare sus límites en $1$ y sus valores en ese punto.
\end{ejemplo}
\ifmostrarSoluciones
\begin{solucion}
Ambas funciones tienen límite $2$ en $1$, porque coinciden cerca de ese punto cuando $x\neq1$. Sin embargo, $g(1)=2$ y $h(1)=7$.
\end{solucion}
\else\vspace{12mm}\fi

\section*{Límites básicos}

\begin{proposicion}[límites fundamentales]
Sean $a,c\in\mathbb{R}$ y $n\in\mathbb{N}^{*}$. Entonces
\[
\lim_{x\to a}c=c,\qquad
\lim_{x\to a}x=a,\qquad
\lim_{x\to a}x^n=a^n.
\]
Cuando $a>0$,
\[
\lim_{x\to a}\sqrt[n]{x}=\sqrt[n]{a}.
\]
Para $n$ par y $a=0$, esta última afirmación se interpreta como límite por la derecha.
\end{proposicion}

\begin{teorema}[álgebra de límites]
Suponga que
\[
  \displaystyle \lim_{x\to a}f(x)=L,
  \qquad
  \displaystyle \lim_{x\to a}g(x)=M,
  \qquad
  \lambda\in\mathbb{R}.
\]
Entonces:
\begin{align*}
\displaystyle \lim_{x\to a}\bigl(f(x)+g(x)\bigr)
  &=\displaystyle \lim_{x\to a}f(x)
    +\displaystyle \lim_{x\to a}g(x)
    =L+M,\\
\displaystyle \lim_{x\to a}\bigl(f(x)-g(x)\bigr)
  &=\displaystyle \lim_{x\to a}f(x)
    -\displaystyle \lim_{x\to a}g(x)
    =L-M,\\
\displaystyle \lim_{x\to a}\bigl(\lambda f(x)\bigr)
  &=\lambda\displaystyle \lim_{x\to a}f(x)
    =\lambda L,\\
\displaystyle \lim_{x\to a}\bigl(f(x)g(x)\bigr)
  &=\left(\displaystyle \lim_{x\to a}f(x)\right)
    \left(\displaystyle \lim_{x\to a}g(x)\right)
    =LM.
\end{align*}
Si $M\neq0$, también
\[
  \displaystyle \lim_{x\to a}\frac{f(x)}{g(x)}
  =\frac{\displaystyle \lim_{x\to a}f(x)}
         {\displaystyle \lim_{x\to a}g(x)}
  =\frac{L}{M}.
\]
Además, si $n\in\mathbb{N}^{*}$, entonces
\[
  \displaystyle \lim_{x\to a}\bigl(f(x)\bigr)^n
  =\left(\displaystyle \lim_{x\to a}f(x)\right)^n
  =L^n.
\]
Si $m\in\mathbb{N}^{*}$ y la expresión radical está definida para los valores de $x$ cercanos a $a$, entonces
\[
  \displaystyle \lim_{x\to a}\sqrt[m]{f(x)}
  =\sqrt[m]{\displaystyle \lim_{x\to a}f(x)}
  =\sqrt[m]{L}.
\]
\end{teorema}

\begin{nota}
La condición $M\neq0$ en la regla del cociente es indispensable. Una forma $0/0$ es \textbf{indeterminada}: distintos cocientes pueden presentar esa forma y tener límites diferentes o no tener límite.
\end{nota}

\section*{Sustitución directa}

\begin{proposicion}[funciones algebraicas]
Sea $a\in\mathbb{R}$.
\begin{enumerate}
  \item Si $p$ es un polinomio, entonces
  \[
    \lim_{x\to a}p(x)=p(a).
  \]
  \item Si $r=p/q$ es racional y $q(a)\neq0$, entonces
  \[
    \lim_{x\to a}r(x)=r(a).
  \]
  \item Para expresiones radicales, la sustitución es válida cuando $a$ y los valores próximos considerados pertenecen al dominio de la raíz.
\end{enumerate}
\end{proposicion}

\begin{ejemplo}[sustitución directa]
Calcule
\[
\lim_{x\to-2}(3x^3-x+4),\qquad
\lim_{x\to1}\frac{x^2+2}{x+3},\qquad
\lim_{x\to5}\sqrt{x+4}.
\]
\end{ejemplo}
\ifmostrarSoluciones
\begin{solucion}
Los denominadores y radicales son admisibles en los puntos indicados. Por sustitución,
\[
-18,\qquad \frac34,\qquad 3,
\]
respectivamente.
\end{solucion}
\else\vspace{15mm}\fi

\section*{Indeterminaciones algebraicas}

\begin{criterio}[simplificación local]
Si la sustitución produce $0/0$:
\begin{enumerate}
  \item factorice numerador y denominador o racionalice la expresión;
  \item simplifique solamente para $x\neq a$;
  \item calcule el límite de la expresión simplificada;
  \item conserve las restricciones del problema original.
\end{enumerate}
La justificación es que ambas expresiones son localmente iguales alrededor de $a$.
\end{criterio}

\begin{ejemplo}[factorización]
Calcule
\[
\lim_{x\to3}\frac{x^2-9}{x-3}.
\]
\end{ejemplo}
\ifmostrarSoluciones
\begin{solucion}
La sustitución produce $0/0$. Para $x\neq3$,
\[
\frac{x^2-9}{x-3}=\frac{(x-3)(x+3)}{x-3}=x+3.
\]
Las funciones son localmente iguales alrededor de $3$; por tanto, el límite es $6$.
\end{solucion}
\else\vspace{16mm}\fi

\begin{ejemplo}[racionalización del numerador]
Calcule
\[
\lim_{x\to0}\frac{\sqrt{x+4}-2}{x}.
\]
\end{ejemplo}
\ifmostrarSoluciones
\begin{solucion}
Se multiplica por el conjugado:
\begin{align*}
\frac{\sqrt{x+4}-2}{x}
&=\frac{(\sqrt{x+4}-2)(\sqrt{x+4}+2)}{x(\sqrt{x+4}+2)}\\
&=\frac{1}{\sqrt{x+4}+2},\qquad x\neq0.
\end{align*}
En consecuencia, el límite es $1/4$.
\end{solucion}
\else\vspace{18mm}\fi

\begin{ejemplo}[racionalización del denominador]
Calcule
\[
\lim_{x\to4}\frac{x-4}{\sqrt{x}-2}.
\]
\end{ejemplo}
\ifmostrarSoluciones
\begin{solucion}
Para $x\neq4$,
\[
\frac{x-4}{\sqrt{x}-2}
=\frac{(\sqrt{x}-2)(\sqrt{x}+2)}{\sqrt{x}-2}
=\sqrt{x}+2.
\]
Luego el límite es $4$.
\end{solucion}
\else\vspace{15mm}\fi

\begin{nota}
No debe cancelarse un sumando. Solo pueden cancelarse factores no nulos. Por ejemplo, en $(x^2-9)/(x-3)$ se factoriza primero y se cancela el factor $x-3$ bajo la condición $x\neq3$.
\end{nota}

\begin{ejercicios}[álgebra y cálculo de límites]
Calcule los siguientes límites e indique la técnica principal.
\begin{enumerate}
  \item $\displaystyle\lim_{x\to2}(4x^3-3x+1)$.
  \item $\displaystyle\lim_{x\to-1}\frac{x^2+3x+5}{x^2+4}$.
  \item $\displaystyle\lim_{x\to4}\sqrt{2x+1}$.
  \item $\displaystyle\lim_{x\to2}\frac{x^2-4}{x-2}$.
  \item $\displaystyle\lim_{x\to-3}\frac{x^2+5x+6}{x+3}$.
  \item $\displaystyle\lim_{x\to1}\frac{x^3-1}{x-1}$.
  \item $\displaystyle\lim_{x\to0}\frac{\sqrt{1+x}-1}{x}$.
  \item $\displaystyle\lim_{x\to5}\frac{x-5}{\sqrt{x-1}-2}$.
  \item $\displaystyle\lim_{x\to2}\frac{\sqrt{x+2}-2}{x-2}$.
  \item $\displaystyle\lim_{x\to0^+}\sqrt{x}$.
\end{enumerate}
\end{ejercicios}

\ifmostrarSoluciones
\begin{solucion}[de los ejercicios]
\begin{enumerate}
  \item $27$, por sustitución directa.
  \item $3/5$, por sustitución directa.
  \item $3$, por sustitución directa.
  \item $4$, pues $(x^2-4)/(x-2)=x+2$ para $x\neq2$.
  \item $-1$, pues $(x+2)(x+3)/(x+3)=x+2$ para $x\neq-3$.
  \item $3$, usando $x^3-1=(x-1)(x^2+x+1)$.
  \item $1/2$, al racionalizar el numerador.
  \item $4$, porque el cociente se simplifica a $\sqrt{x-1}+2$.
  \item $1/4$, al racionalizar el numerador.
  \item $0$; solo se considera la aproximación por la derecha.
\end{enumerate}
\end{solucion}
\fi

\section*{Monotonía del límite y preservación del signo}

En los siguientes resultados, suponemos que existen puntos en el dominio común de las funciones, distintos de $a$, tan próximos a $a$ como se desee. Las desigualdades se exigen para puntos del dominio suficientemente próximos a $a$, con $x\neq a$, y los límites indicados son finitos.

\begin{teorema}[monotonía o comparación de límites]
Si $f(x)\leq g(x)$ cerca de $a$, y
\[
\lim_{x\to a}f(x)=L,\qquad \lim_{x\to a}g(x)=M,
\]
entonces $L\leq M$. En particular, si $m\leq f(x)\leq M$ y el límite de $f$ existe, también pertenece a $[m,M]$.
\end{teorema}

\begin{nota}
Aquí monotonía del límite significa que el paso al límite conserva el orden; no exige que las funciones sean crecientes o decrecientes. Considere las funciones $f,g\colon\mathbb{R}\to\mathbb{R}$ definidas por $f(x)=0$ y $g(x)=x^2$. Para todo $x\neq0$ se cumple la desigualdad estricta
\[
  f(x)<g(x),
  \qquad\text{es decir,}\qquad
  0<x^2.
\]
Sin embargo, los valores positivos de $x^2$ pueden aproximarse tanto como se desee a $0$ cuando $x\to0$. Por ello,
\[
  \lim_{x\to0}f(x)
  =\lim_{x\to0}g(x)
  =0.
\]
Así, una desigualdad estricta entre los valores de dos funciones puede convertirse en una igualdad entre sus límites. En general, de $f(x)<g(x)$ cerca de $a$ solo se puede concluir
\[
  \lim_{x\to a}f(x)
  \leq
  \lim_{x\to a}g(x),
\]
si ambos límites existen; no necesariamente se obtiene una desigualdad estricta.
\end{nota}

\begin{teorema}[preservación del signo]
Si
\[
  \lim_{x\to a}f(x)=L>0,
\]
entonces $f(x)>0$ para todas las entradas suficientemente próximas a $a$, distintas de $a$. Si $L<0$, entonces $f(x)<0$ cerca de $a$.
\end{teorema}

\begin{ejemplo}[signo cerca de un punto]
Sean $p,q\colon\mathbb{R}\to\mathbb{R}$, definidas por $p(x)=x^2-3$ y $q(x)=x-2$. Sin resolver una desigualdad completa, justifique que $p$ es positiva cerca de $2$. ¿Qué sucede con $q$ cerca de $2$?
\end{ejemplo}
\ifmostrarSoluciones
\begin{solucion}
Como
\[
  \lim_{x\to2}p(x)=1>0,
\]
el signo de $p$ es positivo cerca de $2$. En cambio,
\[
  \lim_{x\to2}q(x)=0:
\]
$q$ es negativa a la izquierda y positiva a la derecha. Un límite nulo no garantiza preservación de un signo estricto.
\end{solucion}
\else\vspace{18mm}\fi

\begin{nota}[Observación: material complementario]
El dominio, el recorrido y las gráficas de las funciones seno y coseno se presentan en \emph{Las funciones seno y coseno}, material complementario de esta semana. En los ejemplos siguientes se utilizarán las cotas
\[
  -1\leq\sin(t)\leq1,
  \qquad
  -1\leq\cos(t)\leq1
  \qquad (t\in\mathbb{R}).
\]
\end{nota}

\section*{Teorema del Sanduche}

\begin{teorema}[del Sanduche o del encaje]
Suponga que, para $x$ suficientemente próximo a $a$ y $x\neq a$,
\[
g(x)\leq f(x)\leq h(x).
\]
Si
\[
  \lim_{x\to a}g(x)=\lim_{x\to a}h(x)=L,
\]
entonces
\[
\lim_{x\to a}f(x)=L.
\]
\end{teorema}

\begin{nota}
Las dos cotas deben acercarse al mismo número. El teorema permite establecer la existencia del límite de $f$ sin suponerla previamente. No se necesita que cada factor de una expresión tenga límite.
\end{nota}

\begin{ejemplo}[una oscilación amortiguada]
Calcule
\[
  \lim_{x\to0}x^2\sin(1/x).
\]
\end{ejemplo}
\ifmostrarSoluciones
\begin{solucion}
Para $x\neq0$, $-1\leq\sin(1/x)\leq1$. Al multiplicar por $x^2\geq0$,
\[
-x^2\leq x^2\sin(1/x)\leq x^2.
\]
Las dos cotas tienden a $0$, de modo que el límite es $0$. La oscilación permanece, pero su amplitud disminuye.
\end{solucion}
\else\vspace{18mm}\fi

\begin{center}
\begin{tikzpicture}[x=11.5cm,y=22cm,every node/.style={font=\scriptsize}]
  \draw[->,Gris] (-0.37,0)--(0.38,0) node[right] {$x$};
  \draw[->,Gris] (0,-0.12)--(0,0.125) node[above] {$y$};
  \draw[Rojo,dashed,domain=-0.35:0.35,samples=100] plot (\x,{\x*\x});
  \draw[Rojo,dashed,domain=-0.35:0.35,samples=100] plot (\x,{-\x*\x});
  % Muestreo ilustrativo: las cotas justifican el límite, no el dibujo.
  \draw[Azul,thick,domain=-0.35:-0.012,samples=650] plot (\x,{\x*\x*sin(deg(1/\x))});
  \draw[Azul,thick,domain=0.012:0.35,samples=650] plot (\x,{\x*\x*sin(deg(1/\x))});
  \foreach \x/\etiqueta in {-0.3/$-0.3$,-0.2/$-0.2$,-0.1/$-0.1$,0.1/$0.1$,0.2/$0.2$,0.3/$0.3$}
    \draw (\x,0.002)--(\x,-0.002) node[below] {\etiqueta};
  \node[Rojo,right] at (0.27,0.073) {$y=x^2$};
  \node[Rojo,right] at (0.27,-0.073) {$y=-x^2$};
\end{tikzpicture}
\end{center}
La gráfica muestra el intervalo $-0.35\leq x\leq0.35$. En este acercamiento se observa que las oscilaciones se hacen más rápidas cerca de $0$, mientras quedan encerradas entre $y=-x^2$ y $y=x^2$, una franja cada vez más estrecha.

\begin{ejemplo}[uso de una cota absoluta]
Calcule
\[
  \lim_{x\to0}x\cos(1/x).
\]
\end{ejemplo}
\ifmostrarSoluciones
\begin{solucion}
Como
\[
  |\cos(1/x)|\leq1,
\]
al multiplicar por $|x|\geq0$ obtenemos
\[
  |x\cos(1/x)|
  =|x|\,|\cos(1/x)|
  \leq |x|.
\]
Por tanto,
\[
  -|x|
  \leq x\cos(1/x)
  \leq |x|.
\]
Como
\[
  \lim_{x\to0}-|x|
  =\lim_{x\to0}|x|
  =0,
\]
el teorema del Sanduche permite concluir que
\[
  \lim_{x\to0}x\cos(1/x)=0.
\]
Usar $|x|$ evita invertir incorrectamente las desigualdades cuando $x<0$.
\end{solucion}
\else\vspace{15mm}\fi

\begin{ejercicios}[comparación, signo y encaje]
\begin{enumerate}
  \item Si $2-x^2\leq f(x)\leq2+x^2$ cerca de $0$, determine $\displaystyle\lim_{x\to0}f(x)$.
  \item Calcule $\displaystyle\lim_{x\to0}x^2\cos(3/x)$.
  \item Calcule $\displaystyle\lim_{x\to0}(x-1)x\sin(1/x)$.
  \item Si $|b(x)|\leq5$ cerca de $2$, calcule $\displaystyle\lim_{x\to2}(x-2)b(x)$.
  \item Si $\displaystyle\lim_{x\to a}f(x)=-3$, ¿qué signo tiene $f$ suficientemente cerca de $a$? Justifique.
  \item Si $f(x)>g(x)$ cerca de $a$ y ambos límites existen, ¿se deduce que $\displaystyle\lim_{x\to a}f(x)>\displaystyle\lim_{x\to a}g(x)$? Dé un contraejemplo.
  \item ¿Basta saber que $-1\leq f(x)\leq1$ para afirmar que $\displaystyle\lim_{x\to0}f(x)$ existe?
  \item Si $f(x)\geq4$ cerca de $a$ y su límite existe, ¿puede ese límite ser $3$?
\end{enumerate}
\end{ejercicios}

\ifmostrarSoluciones
\begin{solucion}[de los ejercicios]
\begin{enumerate}
  \item $2$, por el teorema del Sanduche.
  \item $0$, pues $-x^2\leq x^2\cos(3/x)\leq x^2$.
  \item $0$, pues el valor absoluto está acotado por $|x-1||x|\to0$.
  \item $0$, pues el valor absoluto está acotado por $5|x-2|\to0$.
  \item Negativo: las imágenes pueden mantenerse a distancia menor que $3/2$ de $-3$.
  \item No. Por ejemplo, sean $f,g\colon\mathbb{R}\to\mathbb{R}$, definidas por $f(x)=x^2$ y $g(x)=0$. Se cumple $f(x)>g(x)$ para $x\neq0$, pero ambas funciones tienen límite $0$ en $0$.
  \item No. La función $f\colon\mathbb{R}\setminus\{0\}\to\mathbb{R}$, definida por $f(x)=|x|/x$, satisface $-1\leq f(x)\leq1$, pero sus límites laterales en $0$ son distintos.
  \item No; la comparación implica $\displaystyle\lim_{x\to a}f(x)\geq4$.
\end{enumerate}
\end{solucion}
\fi

\fuentes{\textbf{Para ampliar:} \emph{Cálculo en una variable: notas de clase 2018-B}, apartados de límites; J. Stewart, \emph{Cálculo de una variable: trascendentes tempranas}, §§2.2--2.3; y G. Strang--E. Herman, \emph{Cálculo, volumen 1}, OpenStax, §§2.2--2.3.}

\end{document}
